\documentclass[12pt,a4paper]{article}

% ── Paquetes ──────────────────────────────────────────────────────────────────
\usepackage[utf8]{inputenc}
\usepackage[T1]{fontenc}
\usepackage[spanish]{babel}
\usepackage[margin=2cm, top=2.5cm, bottom=2cm]{geometry}
\usepackage{amsmath, amssymb}
\usepackage{booktabs}
\usepackage{array}
\usepackage{tabularx}
\usepackage[table]{xcolor}
\usepackage{enumitem}
\usepackage{fancyhdr}
\usepackage{titlesec}
\usepackage[many, breakable]{tcolorbox}
\usepackage{float}
\usepackage{graphicx}
\usepackage{hyperref}

\hypersetup{colorlinks=true, linkcolor=verdeosc, urlcolor=verdeosc}

% ── Colores institucionales ───────────────────────────────────────────────────
\definecolor{verdeosc}{RGB}{26,107,60}
\definecolor{verdemid}{RGB}{46,139,87}
\definecolor{dorado}{RGB}{212,172,13}
\definecolor{azulclaro}{RGB}{236,240,247}
\definecolor{verdeclaro}{RGB}{232,244,232}
\definecolor{grisclaro}{RGB}{248,248,248}

% ── Encabezados ───────────────────────────────────────────────────────────────
\setlength{\headheight}{30pt}
\pagestyle{fancy}
\fancyhf{}
\fancyhead[L]{\small\color{verdemid}SICVEC 2026}
\fancyhead[R]{\small\color{verdemid}Guía de Presentación}
\fancyfoot[C]{\small\thepage}
\renewcommand{\headrulewidth}{0.6pt}
\renewcommand{\footrulewidth}{0.4pt}
\renewcommand{\headrule}{\hbox to\headwidth{\color{verdeosc}\leaders\hrule height\headrulewidth\hfill}}
\renewcommand{\footrule}{\hbox to\headwidth{\color{dorado}\leaders\hrule height\footrulewidth\hfill}}

% ── Secciones ─────────────────────────────────────────────────────────────────
\titleformat{\section}
  {\normalfont\large\bfseries\color{white}}
  {\fcolorbox{verdeosc}{verdeosc}{\thesection}}{0.7em}{}
  [\vspace{-2pt}\color{verdeosc}\rule{\linewidth}{0.8pt}]

\titleformat{\subsection}
  {\normalfont\normalsize\bfseries\color{verdemid}}{\thesubsection}{0.6em}{}

% ── Cajas ─────────────────────────────────────────────────────────────────────
\tcbset{
  cajainfo/.style={enhanced,breakable,colback=verdeclaro,colframe=verdeosc,
    fonttitle=\bfseries\small\color{white},coltitle=white,
    attach boxed title to top left={yshift=-2mm,xshift=4mm},
    boxed title style={colback=verdeosc},left=6pt,right=6pt,top=10pt,bottom=6pt},
  cajaalerta/.style={enhanced,breakable,colback=yellow!10,colframe=orange,
    left=6pt,right=6pt,top=6pt,bottom=6pt,boxrule=0.8pt}
}

% ============================================================================
\begin{document}
% ============================================================================

\thispagestyle{empty}

\begin{center}
{\Large\bfseries\color{verdeosc} SIMPOSIO INTERNACIONAL DE CIENCIA VERDE Y ECONOMÍA CIRCULAR}\\
{\normalsize\color{verdemid} SICVEC 2026}\\[1cm]
{\Large\bfseries Guía para la Presentación de Resúmenes}\\[0.3cm]
{\normalsize 19 -- 20 de octubre de 2026}\\[1cm]
\end{center}

\section*{Información Importante}

\begin{tcolorbox}[cajainfo, title=Fechas Clave]
\begin{itemize}[label={\textcolor{verdemid}{\textbullet}}]
  \item \textbf{Apertura de convocatoria:} 22 de septiembre de 2026
  \item \textbf{Cierre de recepción de resúmenes:} 4 de octubre de 2026 (23:59)
  \item \textbf{Notificación de aceptación:} 8 -- 9 de octubre de 2026
  \item \textbf{Envío de material final (aceptados):} 14 de octubre de 2026
  \item \textbf{Límite de pago de inscripción:} 15 de octubre de 2026
  \item \textbf{Simposio:} 19 -- 20 de octubre de 2026
\end{itemize}
\end{tcolorbox}

\newpage

% ============================================================================
\section{Modalidades de Presentación}
% ============================================================================

\subsection{1. Ponencia Oral (Oral Presentation)}

Presentación de investigación original con duración de \textbf{20 minutos} (15 min presentación + 5 min preguntas).

\textbf{Requisitos:}
\begin{itemize}[label={\textcolor{verdeosc}{\checkmark}}]
  \item Resumen máximo 300 palabras
  \item Investigación completa con resultados conclusivos
  \item Presentación en PowerPoint
  \item Investigador debe estar disponible para el evento
\end{itemize}

\textbf{Idoneidad:}
\begin{itemize}
  \item Investigaciones finalizadas o en etapa avanzada
  \item Resultados innovadores y pertinentes a SICVEC
  \item Calidad de redacción y claridad conceptual
\end{itemize}

\subsection{2. Poster Científico (Scientific Poster)}

Presentación visual de investigación en formato póster (90 cm alto × 120 cm ancho).

\textbf{Requisitos:}
\begin{itemize}[label={\textcolor{verdeosc}{\checkmark}}]
  \item Resumen máximo 250 palabras
  \item Investigación en curso o proyectos preliminares aceptados
  \item Póster en formato digital (PDF o imagen de alta resolución)
  \item Sesión de networking/posters: 16:00-17:30 (Día 1)
\end{itemize}

\textbf{Idoneidad:}
\begin{itemize}
  \item Proyectos en desarrollo
  \item Resultados preliminares o parciales
  \item Experiencias innovadoras en sostenibilidad
  \item Propuestas de investigación futura
\end{itemize}

% ============================================================================
\section{Estructura de Resúmenes}
% ============================================================================

\subsection{Resumen para Ponencia Oral y Poster}

Los resúmenes deben estructurarse con las siguientes secciones:

\begin{tcolorbox}[cajainfo, title=Estructura estándar]
\begin{enumerate}[label={\textcolor{verdemid}{\textbf{\arabic{*}.}}}]
  \item \textbf{Introducción/Justificación} -- ¿Por qué es importante este tema?
  \item \textbf{Objetivos} -- ¿Qué se pretende lograr?
  \item \textbf{Metodología} -- ¿Cómo se desarrolló la investigación?
  \item \textbf{Resultados/Hallazgos} -- ¿Qué se encontró?
  \item \textbf{Conclusiones/Implicaciones} -- ¿Cuál es el impacto?
\end{enumerate}
\end{tcolorbox}

\subsection{Parámetros de Formato}

\begin{table}[H]
\centering
\small
\begin{tabular}{|l|l|l|}
\hline
\rowcolor{verdeclaro}
\textbf{Parámetro} & \textbf{Ponencia Oral} & \textbf{Poster} \\
\hline
Extensión & Máx. 300 palabras & Máx. 250 palabras \\
\hline
Idioma & Español o Inglés & Español o Inglés \\
\hline
Fuente & Times New Roman, 11pt & Times New Roman, 11pt \\
\hline
Espaciado & 1.5 líneas & 1.5 líneas \\
\hline
Márgenes & 2.5 cm & 2.5 cm \\
\hline
Referencias & Máx. 5 & Máx. 5 \\
\hline
Figuras/Tablas & Máx. 1 & Máx. 2 \\
\hline
\end{tabular}
\end{table}

% ============================================================================
\section{Proceso de Inscripción y Envío}
% ============================================================================

\subsection{Paso 1: Registro de Participante}

Ingresa a la plataforma de inscripción y completa:
\begin{itemize}[label={\textcolor{verdeosc}{\textbullet}}]
  \item Nombre completo y afiliación
  \item Email y teléfono de contacto
  \item Programa académico o profesión
  \item País de procedencia
\end{itemize}

\subsection{Paso 2: Seleccionar Modalidad}

Elige entre:
\begin{enumerate}[label={\textcolor{verdemid}{\arabic{*}.}}]
  \item Ponencia oral
  \item Poster científico
  \item Solo asistencia
\end{enumerate}

\subsection{Paso 3: Enviar Resumen}

Completa el formulario online con:
\begin{itemize}[label={\textcolor{verdeosc}{\textbullet}}]
  \item Título de la investigación
  \item Autores (máx. 5)
  \item Resumen estructurado
  \item Palabras clave (3-5)
  \item Eje temático principal
  \item Resumen pegado como texto en el formulario (sin nombres de autores ni afiliaciones dentro del texto)
\end{itemize}

\subsection{Paso 4: Confirmación}

Recibirás un email de confirmación con:
\begin{itemize}[label={\textcolor{verdemid}{\checkmark}}]
  \item Número de referencia
  \item Estado de envío
  \item Próximos pasos
\end{itemize}

\begin{tcolorbox}[cajaalerta]
\textbf{Importante:} Solo se aceptarán resúmenes enviados hasta las 23:59 del 4 de octubre de 2026 (zona horaria Colombia). Envíos posteriores no serán considerados.
\end{tcolorbox}

% ============================================================================
\section{Evaluación y Peer Review}
% ============================================================================

\subsection{Criterios de Evaluación}

Todos los resúmenes serán evaluados por un comité científico especializado usando una \textbf{rúbrica de 5 criterios}:

\begin{table}[H]
\centering
\small
\begin{tabular}{|l|p{8cm}|c|}
\hline
\rowcolor{verdeclaro}
\textbf{Criterio} & \textbf{Descripción} & \textbf{Peso} \\
\hline
\textbf{Relevancia} & Pertinencia del tema a SICVEC 2026 & 25\% \\
\hline
\textbf{Originalidad} & Novedad y aportación a la ciencia verde & 25\% \\
\hline
\textbf{Calidad Científica} & Rigor metodológico y validez de resultados & 25\% \\
\hline
\textbf{Claridad} & Redacción clara y bien estructurada & 15\% \\
\hline
\textbf{Impacto Potencial} & Aplicabilidad e impacto en sostenibilidad & 10\% \\
\hline
\end{tabular}
\end{table}

\subsection{Proceso de Revisión}

\begin{enumerate}[label={\textcolor{verdemid}{\textbf{Etapa \arabic{*}:}}}]
  \item \textbf{Admisibilidad} (Coordinación Académica) -- Verificación de completitud y formato
  \item \textbf{Asignación de Revisores} -- 2 revisores especialistas por resumen
  \item \textbf{Evaluación de Pares} -- Revisión ciega (double-blind peer review)
  \item \textbf{Recomendación Inicial} -- Revisores proponen aceptación/rechazo
  \item \textbf{Decisión Final} -- Comité científico define resultado
  \item \textbf{Notificación} -- Email con decisión y feedback
\end{enumerate}

\subsection{Categorías de Decisión}

\begin{tcolorbox}[cajainfo, title=Posibles resultados]
\begin{itemize}[label={\textcolor{verdemid}{\textbullet}}]
  \item \textbf{Aceptado:} Será incluido en el programa oficial
  \item \textbf{Aceptado con modificaciones:} Se solicitan cambios menores
  \item \textbf{Rechazado:} No cumple criterios de calidad/relevancia
\end{itemize}
\end{tcolorbox}

% ============================================================================
\section{Derechos de Autor y Memorias}
% ============================================================================

\subsection{Memorias del Evento}

SICVEC 2026 compilará las memorias que incluirán:

\begin{itemize}[label={\textcolor{verdeosc}{\textbullet}}]
  \item Resúmenes de todas las ponencias aceptadas
  \item Posters en formato digital
  \item Listado de participantes
  \item Programa científico
  \item Documentos de política pública generados
\end{itemize}

\subsection{Propiedad Intelectual}

\begin{tcolorbox}[cajainfo, title=Derechos y responsabilidades]
\begin{itemize}[label={\textcolor{verdemid}{\textbullet}}]
  \item Los autores retienen todos los derechos sobre sus trabajos
  \item SICVEC 2026 puede reproducir resúmenes en memorias con citación apropiada
  \item Se requiere consentimiento escrito para cambios posteriores
  \item Los resúmenes no pueden ser duplicados en otros eventos simultáneos
\end{itemize}
\end{tcolorbox}

\subsection{Publicación de Artículos Completos}

Se está coordinando con editoriales para posibilidad de publicación de artículos completos en:

\begin{itemize}[label={\textcolor{verdemid}{\checkmark}}]
  \item Revista indexada en sostenibilidad ambiental
  \item Libro de memoria arbitrado
  \item Foro en línea de SICVEC
\end{itemize}

% ============================================================================
\section{Contacto y Soporte}
% ============================================================================

Para dudas sobre el proceso de envío de resúmenes:

\begin{tcolorbox}[cajainfo, title=Contacto Académico]
\textbf{Coordinadora Académica: María Ximena Díaz}

Email: \texttt{sicvec.academico@unisucre.edu.co}

Teléfono: [Completar]

Horario de atención: Lunes a viernes, 8:00 -- 17:00 (Hora Colombia)
\end{tcolorbox}

\vspace{0.5cm}

\textbf{Sitio web del evento:} [URL del sitio web - Completar]

\textbf{Plataforma de inscripción:} [URL de la plataforma - Completar]

\end{document}
